\documentclass[12pt]{scrartcl}
\usepackage[hagavi]{azzam}

% Modulo dan Keterbagian - Ujian
% 3 soal wajib + 2 soal pilihan (pilih satu untuk full points).
% Sumber tiap soal hanya ditulis di komentar, tidak dimunculkan di PDF.

\title{Modulo dan Keterbagian - Ujian}
\date{12 Oktober 2026}

\begin{document}
\maketitle
\textbf{Aturan umum:}
\begin{itemize}
	\item Tulis \textbf{nama lengkap} dan \textbf{asal sekolah} di pojok kiri atas halaman pertama.
	\item \textbf{Soal bertipe esai}. Sertakan argumentasi atau cara mendapatkan jawaban yang ditanyakan di soal.
	\item Waktu standar untuk mengerjakan semua soal berikut adalah 90--120 menit.
	\item Setiap soal bernilai bilangan bulat antara 0 sampai 10 (inklusif).
	\item Kerjakan \textbf{semua soal wajib} dan \textbf{satu soal pilihan} untuk mendapatkan \textit{full points}.
\end{itemize}

\section{Soal}
\begin{enumerate}

	% Source: OSP SMA 2002 No. 14 (Isian Singkat)
	% Jawaban: 30
	\item (Soal Wajib) Berapakah bilangan bulat positif terbesar yang membagi semua bilangan
	$1^5 - 1$, $2^5 - 2$, \dots, $n^5 - n$?

	% Source: OSP SMA 2002 No. 2 (Esai)
	% Jawaban: p = 3, 5, 9, 35
	\item (Soal Wajib) Tentukan semua bilangan bulat positif $p$ sehingga
	$\frac{3p + 25}{2p - 5}$ juga bulat positif.

	% Source: OSK SMA 2017 No. 3 (Isian Singkat)
	% Jawaban: 8
	\item (Soal Wajib) Banyaknya bilangan asli $k$ yang memenuhi $k \mid (n^7 - n)$
	untuk semua bilangan asli $n$ adalah \dots

	% Source: OSP SMA 2013 No. 3 (Esai)
	% Jawaban: (1) terbukti 2013 valid, (2) semua n genap tidak valid
	\item (Soal Pilihan) Suatu bilangan asli $n$ dikatakan ``valid'' jika
	$1^n + 2^n + 3^n + \dots + m^n$ habis dibagi $1 + 2 + 3 + \dots + m$
	untuk setiap bilangan asli $m$.
	\begin{enumerate}[(i)]
		\item Tunjukkan bahwa $2013$ valid.
		\item Buktikan bahwa ada tak hingga banyak bilangan yang tidak valid.
	\end{enumerate}

	% Source: OSP SMA 2024 No. 5 (Isian Singkat)
	% Jawaban: 176
	\item (Soal Pilihan) Banyaknya bilangan asli $a$ yang kurang dari $209$ sehingga
	$FPB(a, 209) = 1$ dan $a^2 - 1$ bukan kelipatan dari $209$ adalah \dots

\end{enumerate}

\end{document}
